\section{Discussion: extended caveats and stakes}
\label{si:discussion-extras}

\subsection{Smoking--lung-cancer: structural-ceiling derivation and case-control vs cohort}
\label{si:smoking-detail}

Smoking and lung cancer are often cited as the canonical example of a strong, well-established epidemiological causal link. Pooled multiple-adjusted relative risks are about 7 in large cohort meta-analyses \citep{okeeffe2018}. Converting via Equation~\eqref{eq:rr-to-r} at representative population prevalences ($p_X=0.20$ for smoking prevalence, $p_Y=0.06$ for lifetime lung-cancer incidence) gives $r\approx 0.28$, roughly $3\,\sigma_K$ against the NHANES residual scale at $K=10$ ($\sigma_K=0.093$): the smoking--lung-cancer association clears the conventional crud-aware threshold but is, in correlation units, much smaller than the relative-risk scale alone would suggest.

This is in fact close to the structural ceiling on binary correlations at this outcome prevalence:
\begin{equation}
\phi_{\max} = \sqrt{\frac{p_Y(1-p_X)}{p_X(1-p_Y)}} \approx 0.51 \quad \text{here, regardless of RR.}
\end{equation}
Even an arbitrarily large causal effect on a rare binary outcome cannot produce a much larger population correlation in a cohort design. So smoking-cancer at $r\approx 0.28$ is over half the structural maximum, and the test correctly identifies it as exceptional given that ceiling.

Case-control summaries of the same data routinely report $r\approx 0.7$. That figure is a property of the 1:1 case-control sampling design (which artificially sets $p_Y=0.5$), not of the population correlation, and is not the right comparison to a population-cohort $\sigma_K$ such as NHANES (Appendix~\ref{app:scale-bridge}, caveat (vii)). The $\sigma_K$-multiple is design-internal: case-control sampling inflates focal correlations involving the case-defining variable but leaves correlations among other variable pairs (which dominate $\sigma_K$) roughly unchanged, so the same underlying causal relationship can read as $\sim 3\sigma_K$ in cohort units and $\sim 7\sigma_K$ in case-control units.

A second example from a different domain illustrates the upper end without any binary-correlation ceiling: in the HEXACO literature, Dark Triad traits are strongly negatively correlated with Honesty-Humility ($r$s $= -0.72, -0.57, -0.53$); even the smallest of these, $|r|=0.53$, is roughly $3.7\,\sigma_K$ at $K=0$ in HEXACO ($\sigma_0=0.144$; Table~\ref{tab:crud-scale}) \citep{lee2005}. These two textbook examples are illustrative endpoints, chosen to anchor the high end of the magnitude scale, not representative of the typical published effect.

\subsection{Practical stakes: asymmetric costs of skepticism}
\label{si:practical-stakes}

The cost of treating small-effect causal claims skeptically depends on the downstream decisions they inform. Several of those decisions are themselves of modest causal consequence: a regression-discontinuity analysis of NIH R01 funding estimates that receipt of a marginal grant causes fewer than one additional publication over the following five years, with citation effects not statistically distinguishable from zero \citep{jacob2011}. The implication is an asymmetry---raising the bar for small-effect causal claims is unlikely to forgo large substantive gains, because the decisions such claims inform often have only small effects on the outcomes they are meant to improve.

\subsection{Relation to the causal-inference framework}
\label{si:causal-inference-framework}

The crud-aware test does not replace identification arguments. The standard causal-inference framework---directed acyclic graphs, the do-calculus, instruments, regression discontinuities, sensitivity analyses for unobserved confounding \citep{pearl2009,imbens_rubin2015,cinelli2020,oster2019}---is the primary apparatus for translating observational data into causal claims. A small adjusted correlation can be perfectly informative if the DAG identifies the effect; a large one can be uninformative if confounding is severe. Our test is a magnitude diagnostic complementary to identification: it asks ``is the magnitude distinctive against same-pipeline residual dependence,'' which is a different question from ``is the effect identified.'' The test is most useful when identification is uncertain (the typical observational-claim setting), where it adds a magnitude floor on top of an unverified identification argument; it is least useful when identification is clean (e.g., a well-designed RCT), where magnitude alone is already meaningful relative to a sharp null.

\subsection{Sensitivity of \texorpdfstring{$\sigma_K$}{sigma\_K} to discretionary choices}
\label{si:sigma-K-sensitivity}

The reported quantitative multiples ($r/\sigma_K$, etc.) depend on $K$, preprocessing, design type, and reference-dataset choice. The qualitative conclusion---that the residual correlation distribution remains broadly distributed across these choices, and that small adjusted correlations rarely clear the conventional threshold against an approximately matched reference---is robust: $\sigma_K$ at $K=10$ and $K=50$ differ by less than a factor of two in seven of nine datasets (Table~\ref{tab:crud-scale}), and the alt-adjustment results (SI Appendix~\ref{app:alt-adjust}) confirm that Lasso does not collapse the residual distribution either. The quantitative pass/fail at the conventional threshold for a specific paper, however, is sensitive enough that we recommend reporting conclusions across a range of plausible $\sigma_K$ values rather than a single number (SI Appendix~\ref{si:how-to-use}, SI Appendix~\ref{si:reference-catalog}).

\subsection{Why the random-pair null is the right reference, even for ``mechanism-driven'' hypotheses}
\label{si:random-pair-null}

An applied investigator may object that hypothesizing from mechanism (folate$\to$one-carbon metabolism$\to$methylation, smoking$\to$DNA adducts$\to$lung carcinoma) is not the same as drawing a random pair, so the random-pair null is the wrong reference. In narrow settings the objection has force---a single pair pre-specified before any data was looked at, tested in an independent dataset. In practice, however, ``I had a mechanism'' is often a post-hoc rationalization of one story among many that could have been told. The literature contains thousands of plausible mechanistic narratives; selecting which one to test, in which dataset, with which adjustment set, is itself a form of multiple selection. The applied user typically samples mechanisms from a vast latent pool and reports the ones that survive; the magnitude distribution of the reported subset is much closer to the random-pair magnitude distribution than the user's introspective sense of being mechanistically pre-specified would suggest. The crud-aware percentile rank treats this honestly: it asks whether the reported magnitude is what one would see by drawing pairs at random from the domain. If the investigator's mechanism is genuinely privileged, the answer should distinguish it from the bulk; if the mechanism is one of many such stories, the random-pair null is the appropriate reference. The test does not adjudicate which case the investigator is in---that is a question of design, mechanism, and replication---but the random-pair reference is defensible precisely because the alternative (``my mechanism is privileged, so I deserve a narrower null'') is rarely justified in practice.

\subsection{Magnitude-as-descriptor fallback for fields without design leverage}
\label{si:fallback-advice}

For fields where design leverage is structurally unavailable---nutritional, environmental, and lifestyle epidemiology, where randomization of long-term exposure is infeasible and instruments and discontinuities are vanishingly rare---the honest fallback is: report the magnitude rank as a descriptive feature of the result (``$r=0.02$, crud-aware percentile $0.19$''), do not claim causation from magnitude alone, and concentrate causal-interpretation language on the converging-evidence chain (Mendelian randomization where applicable, negative controls, dose-response, biological plausibility) rather than on the magnitude. When neither design leverage nor a converging-evidence chain is available and the magnitude is in the bulk of the empirical null, the appropriate reporting standard is to label the association as not distinctive against background and treat any causal interpretation as provisional.
